\chapter{The Mathieu Group \texorpdfstring{$M_{24}$}{M24}}
\label{ch:m24}

The elliptic genus of a K3 surface, computed in \cref{ch:ellgenus}, is a modular object built
from the geometry of a four-manifold and the representation theory of the $N=4$
superconformal algebra. Its expansion coefficients are integers, and the first of them are
$90,\,462,\,1440,\,4554,\,11592,\dots$ Eguchi, Ooguri and Tachikawa noticed that these are twice
$45,\,231,\,770,\,2277,\,5796$, and that every one of the latter is the dimension of an
irreducible representation of a finite group of order $244\,823\,040$, the largest Mathieu
group $\Mtf$ \source{EOT §1, ll.~197--203, 228--239}. \Cref{ch:moonshine} is about that
observation. Before we can discuss it we must know the group: where it comes from, how big it
is, what its subgroups and representations look like, and what precisely the numbers
$45, 231, 770, \dots$ are. The question of this chapter is therefore: \emph{what is $\Mtf$,
and which of the finite, explicit facts about it has this project handed to Lean?}

A student should care for two reasons. First, $\Mtf$ is not an exotic object bolted onto
string theory from outside: it is the symmetry group of a $24$-dimensional even self-dual
lattice, and \cref{ch:lattices,ch:k3lattice,ch:mukai} have already shown that the
cohomology of K3 is a $24$-dimensional even self-dual lattice. Mukai's theorem, quoted in
\cref{sec:m24:k3}, makes the link precise: every finite group of symplectic automorphisms of
a K3 surface is a subgroup of $\Mtf$. Secondly, this chapter is a case study in how a table
of numbers is handed to a proof assistant \emph{and then checked}. The project's table of
representation dimensions of $\Mtf$ was wrong for months, and a one-line theorem, Burnside's
$\sum\dim^2=|G|$, exposed it (\cref{sec:m24:story}). That story is worth more than the table.

\emph{Conventions.} Groups act on the left; $|G|$ is the order of $G$; $\mathbb F_2=\{0,1\}$
is the field with two elements; the \emph{weight} $\mathrm{wt}(c)$ of a vector
$c\in\mathbb F_2^n$ is the number of its nonzero coordinates. The root lattice $A_1$ is the
rank-one lattice whose generator has $x^2=2$, so $A_1\cong\sqrt2\,\Z\subset\R$ with the
Euclidean form, and its dual $A_1^*=\tfrac1{\sqrt2}\Z$ has generator of norm $\tfrac12$
\source{EOT App.~B, ll.~959--965}. Dimensions of representations are over $\C$.

%=====================================================================================
\section{Finite simple groups in one page}\label{sec:m24:simple}

A group $G$ is \emph{simple}\index{simple group} if it has no normal subgroups other than
$\{e\}$ and $G$. Simple groups are to finite groups what primes are to integers: every finite
group has a composition series whose successive quotients are simple, and these quotients are
determined by the group (Jordan--H\"older). The classification of finite simple groups
(standard; not checked against a source in this book's library) says that a finite simple
group is cyclic of prime order, alternating $A_n$ with $n\ge5$, a group of Lie type over a
finite field such as $PSL(n,q)$, or one of $26$ \emph{sporadic}\index{sporadic group} groups
that belong to no infinite family. The five Mathieu groups $M_{11}, M_{12}, M_{22}, M_{23},
M_{24}$ are the smallest and oldest sporadic groups; the largest sporadic group is the
Monster, whose representation dimensions $1,\,196883,\dots$ appear in the coefficients of the
modular function $J(q)=q^{-1}+196884\,q+\dots$ (monstrous moonshine, \source{EOT §1,
ll.~244--256}).

Everything in this chapter concerns $\Mtf$\index{M24@$M_{24}$}, and the first fact to record is its order
\source{EOT App.~A, ll.~331--334}:
\begin{equation}\label{eq:m24:order}
 |\Mtf| \;=\; 2^{10}\cdot3^3\cdot5\cdot7\cdot11\cdot23 \;=\; 244\,823\,040 .
\end{equation}
The right-hand side is a number one can check by hand ($2^{10}=1024$, $1024\cdot27=27648$,
$27648\cdot35=967680$, $967680\cdot11=10\,644\,480$, $10\,644\,480\cdot23=244\,823\,040$), and
it is the first of the arithmetic facts that the project has kernel-checked
(\lean{M24\_order}, \cref{sec:m24:lean}). The prime factorization already constrains the
group a great deal: by Lagrange's theorem\index{Lagrange's theorem} the order of every
subgroup and of every element divides $|\Mtf|$, so no element has order $13$, $25$ or $49$
(whereas orders $9$ and $16$ are allowed by Lagrange, though we shall see in
\cref{sec:m24:chars} that they do not occur), and by Sylow's theorem a subgroup of order
$23$ exists (and is cyclic, since $23$ is prime).
Note that $23^2\nmid|\Mtf|$, and that the primes $2,3,5,7,11,23$ are exactly the primes EOT
observe to recur in the factorizations of the moonshine coefficients $A_n$ \source{EOT §1,
footnote~2, ll.~260--262}.

\begin{pitfallbox}[An order is not a group]
Knowing $|G|$ does not determine $G$: there are many groups of order $244\,823\,040$ (the
cyclic group is one of them). Every Lean theorem quoted in this chapter that mentions
$244\,823\,040$ is a statement about that \emph{numeral}. That the numeral is the order of a
simple group called $\Mtf$ is Tier~L, quoted from EOT, and nothing about a group is
constructed in the project's libraries. Read the statements, not the names.
\end{pitfallbox}

%=====================================================================================
\section{Binary codes and the extended Golay code}\label{sec:m24:golay}

Of the many equivalent definitions of $\Mtf$, EOT choose the one closest to string theory:
$\Mtf$ is the symmetry group of a particular even self-dual lattice of rank $24$ built on
$A_1^{24}$, and that lattice is encoded by a binary code \source{EOT App.~B, ll.~959--982}.
We reconstruct their argument in full, because every step is a two-line computation and
because the same construction with $n=8$ produces $E_8$, which the student already knows.

\begin{definition}[Binary linear code]\label{def:m24:code}
A \emph{binary linear code}\index{binary code} of length $n$ and dimension $k$ is a
$k$-dimensional $\mathbb F_2$-subspace $C\subseteq\mathbb F_2^n$. Its \emph{minimum
distance}\index{minimum distance} $d$ is the smallest weight of a nonzero codeword, and one
writes $C$ as an $[n,k,d]$ code. The \emph{dual code} is
$C^\perp=\{x\in\mathbb F_2^n:\ x\cdot c=0\ \text{for all }c\in C\}$ with the dot product
$x\cdot c=\sum_ix_ic_i\bmod2$; $C$ is \emph{self-dual} if $C=C^\perp$ and \emph{doubly even}
if every codeword has weight divisible by $4$.
\end{definition}

Since $\dim C+\dim C^\perp=n$, a self-dual code has $k=n/2$ and $n$ is even. A code with
$2^k$ words can correct any $t=\lfloor(d-1)/2\rfloor$ errors: two balls of radius $t$ around
distinct codewords are disjoint because the codewords differ in at least $d\ge2t+1$ places.

\subsection*{From a code to a lattice}
Consider the lattice $A_1^n=(\sqrt2\,\Z)^n\subset\R^n$ with the Euclidean form. It is even
(every vector has $x^2=2\sum m_i^2$) but far from unimodular: its Gram matrix is $2\cdot\mathbf1_n$,
of determinant $2^n$, and its dual is $(A_1^*)^n=(\tfrac1{\sqrt2}\Z)^n$. Any lattice $N$ with
\begin{equation}\label{eq:m24:sandwich}
 A_1^n\ \subset\ N\ \subset\ (A_1^*)^n
\end{equation}
is determined by the subgroup $C:=N/A_1^n$ of $(A_1^*)^n/A_1^n\cong(\Z/2)^n=\mathbb F_2^n$,
i.e.\ by a binary code \source{EOT App.~B, ll.~966--970}. Concretely, every
$x\in(A_1^*)^n$ can be written uniquely as
\begin{equation}\label{eq:m24:xcm}
 x=\tfrac1{\sqrt2}\,(c+2m),\qquad c\in\{0,1\}^n,\ m\in\Z^n ,
\end{equation}
and $x\in N$ iff $c\bmod2\in C$. Now compute the three lattice properties in terms of $C$.

\emph{Integrality.} For $x=\tfrac1{\sqrt2}(c+2m)$ and $y=\tfrac1{\sqrt2}(c'+2m')$,
\begin{equation}\label{eq:m24:pairing}
 x\cdot y=\tfrac12\sum_i(c_i+2m_i)(c'_i+2m'_i)=\tfrac12\,(c\cdot c')_{\Z}+\text{(integer)},
\end{equation}
where $(c\cdot c')_\Z=\sum_ic_ic'_i$ is the integer dot product. So $N$ is integral iff
$c\cdot c'\equiv0\pmod2$ for all $c,c'\in C$, i.e.\ iff $C\subseteq C^\perp$.

\emph{Evenness.} Setting $y=x$,
\begin{equation}\label{eq:m24:norm}
 x^2=\tfrac12\sum_i(c_i+2m_i)^2=\tfrac12\,\mathrm{wt}(c)+2\sum_i(c_im_i+m_i^2),
\end{equation}
using $c_i^2=c_i$. Hence $x^2\in2\Z$ for all $x\in N$ iff $\mathrm{wt}(c)\equiv0\pmod4$ for
all $c\in C$: \emph{$N$ is even iff $C$ is doubly even} \source{EOT App.~B, ll.~972--973}.

\emph{Unimodularity.} The index $[N:A_1^n]=|C|=2^k$, and for any finite-index overlattice
$\operatorname{disc}N=\operatorname{disc}(A_1^n)/[N:A_1^n]^2=2^n/2^{2k}$ (\cref{ch:lattices}).
So $N$ is unimodular iff $k=n/2$; together with integrality this says $C=C^\perp$:
\emph{$N$ is self-dual iff $C$ is self-dual} \source{EOT App.~B, ll.~971--972}.

\emph{Roots.} A vector of $N$ with $x^2=2$ must, by \eqref{eq:m24:norm}, have
$\mathrm{wt}(c)\in\{0,4\}$: if $c=0$ then $\sum m_i^2=1$ and $x=\pm\sqrt2\,e_i$ is one of the
$2n$ roots of $A_1^n$; if $\mathrm{wt}(c)=4$ then $\sum_i(c_im_i+m_i^2)=0$ forces $m_i=0$ off
the support of $c$ and $m_i\in\{0,-1\}$ on it, giving $2^4=16$ vectors
$\tfrac1{\sqrt2}(\pm1,\pm1,\pm1,\pm1,0,\dots)$ for each weight-$4$ word. Weight $\ge8$ gives
$x^2\ge4$. Therefore
\begin{equation}\label{eq:m24:roots}
 \#\{x\in N:\ x^2=2\}\;=\;2n+16\,N_4(C),\qquad N_4(C):=\#\{c\in C:\ \mathrm{wt}(c)=4\}.
\end{equation}

\begin{example}[The Hamming code and $E_8$]\label{exa:m24:hamming}\index{Hamming code}\index{Construction A}
Take $n=8$ and let $C_8\subset\mathbb F_2^8$ be spanned by the rows of
\[
 G_8=\begin{pmatrix}1&1&1&1&0&0&0&0\\0&0&1&1&1&1&0&0\\0&0&0&0&1&1&1&1\\0&1&0&1&0&1&0&1\end{pmatrix}.
\]
The rows are pairwise orthogonal (each pair overlaps in $0$ or $2$ positions) and have weight
$4$, so $C_8\subseteq C_8^\perp$, and since $\dim C_8=4=8/2$ it is self-dual. Listing the
$16$ codewords gives the weight distribution $1$ word of weight $0$, $14$ of weight $4$, $1$
of weight $8$ (a symbolic check; the reader should do it by hand once). So $C_8$ is a doubly
even self-dual $[8,4,4]$ code (the extended Hamming code), and by
\eqref{eq:m24:roots} the lattice $N_8$ is even, unimodular, of rank $8$, with
\[
 2\cdot8+16\cdot14=16+224=240
\]
roots. By the classification of \cref{ch:lattices}, $N_8\cong E_8$, and $240$ is indeed the
number of roots of $E_8$ found there. This is a satisfying check that the bookkeeping
\eqref{eq:m24:norm}--\eqref{eq:m24:roots} is right.
\end{example}

\subsection*{The Golay code}
Now take $n=24$. EOT ask for an even self-dual $N$ in \eqref{eq:m24:sandwich} \emph{whose only
roots are the $48$ roots of $A_1^{24}$}, i.e.\ $N_4(C)=0$ \source{EOT App.~B, ll.~973--977}.
By the three computations above, $C$ must be a doubly even self-dual $[24,12,d]$ code with
$d\ge8$. Such a code exists and is unique up to permutation of coordinates
\source{EOT App.~B, ll.~977--978}; this uniqueness is a theorem of coding theory that we take
from EOT as Tier~L. It is the \emph{extended binary Golay code}\index{Golay code}
$\mathcal G_{24}$, a $[24,12,8]$ code.

\begin{definition}[$\Mtf$]\label{def:m24:M24}
$\Mtf$ is the subgroup of the symmetric group $S_{24}$, acting on $\mathbb F_2^{24}$ by
permuting coordinates, that maps $\mathcal G_{24}$ to itself: $\Mtf=\operatorname{Aut}(\mathcal
G_{24})$ \source{EOT App.~B, ll.~978--979}.
\end{definition}

Two remarks make this concrete. First, the weight distribution of $\mathcal G_{24}$ is
\begin{equation}\label{eq:m24:weights}
\begin{aligned}
 &\text{weight}&&0&&8&&12&&16&&24\\
 &\text{number of words}&&1&&759&&2576&&759&&1
\end{aligned}
\qquad\qquad 1+759+2576+759+1=4096=2^{12}.
\end{equation}
This is a symbolic check: the author built the code as the extended quadratic-residue code
of length $23$ (the cyclic code spanned by the shifts of the indicator vector of the
nonzero squares modulo $23$, extended by a parity bit), verified $\dim=12$, and counted
weights with a short script. The interpretation of $759$ as the number of \emph{octads} and
the fact that any $5$ of the $24$ points lie in exactly one octad, $759=\binom{24}5/\binom85$,
are standard (Steiner system\index{Steiner system} $S(5,8,24)$) and not checked against a
source in this book's library. Secondly, the Golay code corrects $t=\lfloor(8-1)/2\rfloor=3$ errors, and its
codewords number $2^{12}=4096$; these two numerals are the whole content of the project's
``Golay'' theorems (\cref{sec:m24:lean}).

\begin{physicsbox}[Why a string theorist likes this definition]
The lattice $N=N_{24}$ built from $\mathcal G_{24}$ is even, self-dual and positive
definite of rank $24$, so it is one of the $24$ Niemeier lattices\index{Niemeier lattice} of
\cref{ch:lattices};
it defines a chiral conformal field theory with $c=24$ whose current algebra is
$A_1^{24}$ (the $48$ roots are the $48$ charged currents). $\Mtf$ permutes the $24$ copies
of $A_1$ and preserves the lattice, so it is a discrete symmetry of that chiral CFT
\source{EOT App.~B, ll.~979--982}. The K3 cohomology lattice $H^*(K3,\Z)$ is also even
and self-dual of rank $24$, but of signature $(4,20)$ (\cref{ch:mukai}); EOT locate the
whole K3--$\Mtf$ connection in this coincidence of rank and type \source{EOT App.~B,
ll.~984--990}.
\end{physicsbox}

\begin{figure}[t]\centering
\begin{tikzpicture}[>=Latex,node distance=9mm]
\node (A) {$A_1^{24}=(\sqrt2\,\Z)^{24}$};
\node (N) [right=22mm of A] {$N$};
\node (D) [right=22mm of N] {$(A_1^*)^{24}=(\tfrac1{\sqrt2}\Z)^{24}$};
\draw[->] (A) -- node[above]{\small index $2^{12}$} (N);
\draw[->] (N) -- node[above]{\small index $2^{12}$} (D);
\node (F) [below=9mm of N] {$\mathcal G_{24}=N/A_1^{24}\ \subset\ \mathbb F_2^{24}=(A_1^*)^{24}/A_1^{24}$};
\draw[->,dashed] (N) -- (F);
\node (M) [below=8mm of F] {$\Mtf=\operatorname{Aut}(\mathcal G_{24})\subset S_{24}$ permutes the $24$ coordinates};
\end{tikzpicture}
\caption{The lattice sandwich \eqref{eq:m24:sandwich}. $N$ is even iff the code is doubly
even, self-dual iff the code is self-dual, and has no roots beyond those of $A_1^{24}$ iff
the code has no words of weight $4$. The total index $[(A_1^*)^{24}:A_1^{24}]=2^{24}$ splits
as $2^{12}\cdot2^{12}$ exactly when $N$ is unimodular.}
\label{fig:m24:sandwich}
\end{figure}

%=====================================================================================
\section{Transitivity and the stabilizer tower}\label{sec:m24:tower}

\begin{definition}[$k$-transitive action]\label{def:m24:transitive}
A group $G$ acting on a finite set $\Omega$ is \emph{$k$-transitive}\index{transitive
action} if for any two ordered $k$-tuples of distinct points $(a_1,\dots,a_k)$,
$(b_1,\dots,b_k)$ there is $g\in G$ with $g a_i=b_i$ for all $i$. The \emph{stabilizer} of a
point $a$ is $G_a=\{g:\ ga=a\}$, and the orbit--stabilizer theorem\index{orbit--stabilizer
theorem} gives $|G|=|G\cdot a|\cdot|G_a|$.
\end{definition}

If $G$ is transitive on $\Omega$ with $|\Omega|=n$, then $|G|=n\,|G_a|$, and $G$ is
$k$-transitive iff $G_a$ is $(k-1)$-transitive on $\Omega\setminus\{a\}$. Iterating,
a $k$-transitive group of degree $n$ satisfies
\begin{equation}\label{eq:m24:ktrans}
 |G|=n(n-1)\cdots(n-k+1)\cdot|G_{a_1\dots a_k}| ,
\end{equation}
where $G_{a_1\dots a_k}$ is the pointwise stabilizer of $k$ points. The Mathieu group $\Mtf$
acts on the $24$ coordinates of the Golay code, and this action is $5$-transitive (standard;
the project's \lean{MathieuTower.lean} docstring states it, but it is not checked against a
source in this book's library, and it is not formalized). Successively fixing points defines
the \emph{stabilizer tower}\index{Mathieu tower}
\begin{equation}\label{eq:m24:tower}
 \Mtf\ \supset\ M_{23}\ \supset\ M_{22}\ \supset\ M_{21}\cong PSL(3,4),
\end{equation}
where $M_{23}$ is the stabilizer of one point, $M_{22}$ of two, $M_{21}$ of three. Each step
is an instance of \eqref{eq:m24:ktrans}: since $\Mtf$ is transitive on $24$ points,
$M_{23}$ on the remaining $23$, and $M_{22}$ on the remaining $22$,
\begin{equation}\label{eq:m24:indices}
 |\Mtf|=24\,|M_{23}|,\qquad |M_{23}|=23\,|M_{22}|,\qquad |M_{22}|=22\,|M_{21}| .
\end{equation}
Starting from \eqref{eq:m24:order} and dividing,
\begin{align}
 |M_{23}|&=244\,823\,040/24=10\,200\,960=2^7\cdot3^2\cdot5\cdot7\cdot11\cdot23,\label{eq:m24:M23}\\
 |M_{22}|&=10\,200\,960/23=443\,520=2^7\cdot3^2\cdot5\cdot7\cdot11,\label{eq:m24:M22}\\
 |M_{21}|&=443\,520/22=20\,160=2^6\cdot3^2\cdot5\cdot7.\label{eq:m24:M21}
\end{align}
The last number is recognizable. The group $GL(3,q)$ of invertible $3\times3$ matrices over
the field with $q$ elements has order $(q^3-1)(q^3-q)(q^3-q^2)$ (choose the three columns
to be linearly independent: $q^3-1$ choices for the first, $q^3-q$ for the second, $q^3-q^2$
for the third); dividing by the $q-1$ nonzero determinants gives $|SL(3,q)|$, and by the
$\gcd(3,q-1)$ scalar matrices of determinant $1$ gives $|PSL(3,q)|$. For $q=4$:
\begin{equation}\label{eq:m24:psl34}
 |PSL(3,4)|=\frac{(64-1)(64-4)(64-16)}{(4-1)\cdot\gcd(3,3)}=\frac{63\cdot60\cdot48}{3\cdot3}
 =\frac{181\,440}{9}=20\,160,
\end{equation}
in agreement with \eqref{eq:m24:M21}. That $M_{21}$ is in fact isomorphic to $PSL(3,4)$ is
Tier~L (standard; stated in the project's docstring, not in this book's library); that the
\emph{orders} agree is arithmetic, and it is what the project proves.

Going two steps further down the tower, $5$-transitivity and \eqref{eq:m24:ktrans} give
\begin{equation}\label{eq:m24:fivetrans}
 24\cdot23\cdot22\cdot21\cdot20=5\,100\,480,\qquad
 \frac{|\Mtf|}{5\,100\,480}=48 ,
\end{equation}
so the pointwise stabilizer of five points has order $48$ (and $\Mtf$ is not \emph{sharply}
$5$-transitive: the map from $\Mtf$ to ordered $5$-tuples is $48$-to-one). Read the other
way, \eqref{eq:m24:fivetrans} is a divisibility test: any $5$-transitive group of degree $24$
has order divisible by $5\,100\,480=2^7\cdot3^2\cdot5\cdot7\cdot11\cdot23$, and
\eqref{eq:m24:order} passes.

\begin{figure}[t]\centering
\begin{tikzpicture}[>=Latex,every node/.style={font=\small}]
\node[draw,rounded corners,minimum width=34mm] (m24) at (0,0) {$\Mtf$, order $244\,823\,040$};
\node[draw,rounded corners,minimum width=34mm] (m23) at (0,-1.3) {$M_{23}$, order $10\,200\,960$};
\node[draw,rounded corners,minimum width=34mm] (m22) at (0,-2.6) {$M_{22}$, order $443\,520$};
\node[draw,rounded corners,minimum width=34mm] (m21) at (0,-3.9) {$M_{21}\cong PSL(3,4)$, order $20\,160$};
\draw[->] (m24) -- node[right]{fix $1$ point, index $24$} (m23);
\draw[->] (m23) -- node[right]{fix a $2$nd point, index $23$} (m22);
\draw[->] (m22) -- node[right]{fix a $3$rd point, index $22$} (m21);
\node[left=6mm of m24,align=right] {$5$-transitive\\on $24$ points};
\node[left=6mm of m23,align=right] {$4$-transitive\\on $23$ points};
\node[left=6mm of m22,align=right] {$3$-transitive\\on $22$ points};
\node[left=6mm of m21,align=right] {$2$-transitive\\on $21$ points};
\end{tikzpicture}
\caption{The point-stabilizer tower \eqref{eq:m24:tower}. Each arrow divides the order by the
number of points the group above it acts on transitively; the product
$24\cdot23\cdot22\cdot20\,160$ returns $|\Mtf|$ (\lean{mathieu\_tower\_consistent}).}
\label{fig:m24:tower}
\end{figure}

\begin{pitfallbox}[Index relations do not prove existence]
Equations \eqref{eq:m24:indices} are consequences of transitivity, not evidence for it.
Four numerals $a=24b$, $b=23c$, $c=22d$ are consistent with infinitely many groups (and with
none). The Lean theorems \lean{M24\_stabilizer\_index}, \lean{M23\_stabilizer\_index},
\lean{M22\_stabilizer\_index} of \cref{sec:m24:lean} verify exactly these numeral equations
for the ATLAS values of $|M_{23}|,|M_{22}|,|M_{21}|$; they say nothing about groups acting on
sets. The ATLAS values themselves are Tier~L and are not in this book's library of sources.
\end{pitfallbox}

%=====================================================================================
\section{Conjugacy classes and irreducible representations}\label{sec:m24:chars}

\subsection*{Classes}
Two elements $g,h$ of a group are \emph{conjugate} if $h=xgx^{-1}$ for some $x$; conjugacy
classes\index{conjugacy class} partition the group, and a class function (such as the
character of a representation) is constant on each. $\Mtf$ has $26$ conjugacy classes
\source{EOT App.~A, l.~336}, named in the ATLAS convention by the order of their elements
and a letter \source{GHV Table~1, ll.~1916--1922, 1928--1930}:
\begin{equation}\label{eq:m24:classes}
\begin{aligned}
 &1A,\ 2A,\ 2B,\ 3A,\ 3B,\ 4A,\ 4B,\ 4C,\ 5A,\ 6A,\ 6B,\ 7A,\ 7B,\ 8A,\ 10A,\ 11A,\\
 &12A,\ 12B,\ 14A,\ 14B,\ 15A,\ 15B,\ 21A,\ 21B,\ 23A,\ 23B .
\end{aligned}
\end{equation}
The element orders that occur are thus $1,2,3,4,5,6,7,8,10,11,12,14,15,21,23$, every one a
divisor of \eqref{eq:m24:order} as Lagrange requires (symbolic check). Two features will
matter in \cref{ch:moonshine}. First, if $\gcd(a,o(g))=1$ then $g^a$ lies in the same class
as $g$ \source{GHV §3.1, ll.~694--697}: the character table\index{character table} of
$\Mtf$ is ``almost rational'',
its only irrational entries being of the form $e_p^\pm=(-1\pm\sqrt{-p})/2$ \source{EOT
App.~A, eq.~(A.2), ll.~339--343}, which is why the classes come in lettered pairs such as
$7A/7B$ or $23A/23B$. Secondly, the classes $1A,2A,3A,4B,5A,6A$ contain elements of $M_{23}$
(they fix a point), whereas $2B,3B,4A,4C$ do not; the latter are exactly the classes on
which the permutation character $\mathrm{Tr}_{23\oplus1}(g)$ vanishes \source{GHV §3.1,
ll.~808--813}. The number of fixed points of $g$ on the $24$ coordinates is the character of
the $24$-dimensional permutation representation, which splits as $\mathbf{24}=\mathbf1\oplus\mathbf{23}$
\source{GHV §3, l.~555}.

\subsection*{Irreducible representations and their dimensions}
A finite group has as many irreducible complex representations as conjugacy classes, so
$\Mtf$ has $26$ of them \source{EOT App.~A, l.~336}. Their dimensions, in increasing order,
are \source{EOT App.~A, eq.~(A.3), ll.~345--351}
\begin{equation}\label{eq:m24:dims}
\begin{aligned}
 &1,\ 23,\ 45,\ 45,\ 231,\ 231,\ 252,\ 253,\ 483,\ 770,\ 770,\ 990,\ 990,\ 1035,\ 1035,\ 1035,\\
 &1265,\ 1771,\ 2024,\ 2277,\ 3312,\ 3520,\ 5313,\ 5796,\ 5544,\ 10395 .
\end{aligned}
\end{equation}
(EOT list $5796$ before $5544$; we keep their order because the project's Lean table does.)
The representations of dimensions $45,231,770,990,1035$ come in complex-conjugate pairs,
and there is a third, real, representation of dimension $1035$ \source{EOT App.~A,
eq.~(A.4), ll.~353--358}; so $10$ of the $26$ representations are complex and $16$ are real.
The repeated entries in \eqref{eq:m24:dims} are therefore not typographical: $45$ and $\overline{45}$
are different representations of the same dimension, and a table that lists $45$ once is
\emph{wrong}. This is precisely the kind of error that \cref{sec:m24:story} is about.

Several of these numbers have a direct meaning on the $24$ points, which is a useful way to
remember them (the interpretations are standard and not checked against a source in this
book's library; the arithmetic is a symbolic check): $23$ is the nontrivial part of the
permutation representation; $\binom{24}2=276=1+23+252$ is the dimension of $\Lambda^2$ of
the permutation representation; $253=759\cdot8/24$ is the number of octads through a
point; $2024=\binom{24}3$; and $1771=\binom{24}4/6$.

\subsection*{Burnside's sum of squares}
\begin{proposition}[Burnside]\label{prop:m24:burnside}\index{Burnside's sum of squares}
For a finite group $G$ with irreducible complex representations $V_1,\dots,V_r$,
\begin{equation}\label{eq:m24:burnside}
 \sum_{i=1}^r(\dim V_i)^2=|G| .
\end{equation}
\end{proposition}
\begin{proof}[Derivation]
Let $G$ act on the vector space $\C[G]$ of formal combinations of group elements by left
multiplication (the \emph{regular representation}\index{regular representation}). Its
character is $\chi_{\mathrm{reg}}(e)=|G|$ and $\chi_{\mathrm{reg}}(g)=0$ for $g\ne e$, because
$g\ne e$ fixes no basis vector. The multiplicity of $V_i$ in $\C[G]$ is the inner product
$\langle\chi_{\mathrm{reg}},\chi_i\rangle=\frac1{|G|}\sum_g\chi_{\mathrm{reg}}(g)\overline{\chi_i(g)}
=\frac1{|G|}\cdot|G|\cdot\dim V_i=\dim V_i$. Comparing dimensions of $\C[G]=\bigoplus_iV_i^{\oplus\dim V_i}$
gives \eqref{eq:m24:burnside}. (Standard; uses orthogonality of characters, which is not
proved in this book.)
\end{proof}

\begin{example}[Burnside for $\Mtf$, with numbers]\label{exa:m24:burnside}
Squaring and adding the $26$ entries of \eqref{eq:m24:dims} in blocks:
\begin{align*}
 1^2+23^2&=530,\\
 2\cdot45^2+2\cdot231^2&=4\,050+106\,722=110\,772,\\
 252^2+253^2+483^2&=63\,504+64\,009+233\,289=360\,802,\\
 2\cdot770^2+2\cdot990^2+3\cdot1035^2&=1\,185\,800+1\,960\,200+3\,213\,675=6\,359\,675,\\
 1265^2+1771^2+2024^2+2277^2&=1\,600\,225+3\,136\,441+4\,096\,576+5\,184\,729=14\,017\,971,\\
 3312^2+3520^2+5313^2&=10\,969\,344+12\,390\,400+28\,227\,969=51\,587\,713,\\
 5796^2+5544^2+10395^2&=33\,593\,616+30\,735\,936+108\,056\,025=172\,385\,577 .
\end{align*}
The total is $530+110\,772+360\,802+6\,359\,675+14\,017\,971+51\,587\,713+172\,385\,577
=244\,823\,040=|\Mtf|$, as \eqref{eq:m24:burnside} demands. This is the computation that Lean
performs in \lean{M24RepDim\_sum\_sq} (\cref{sec:m24:lean}). A second, independent sanity
check is that every entry of \eqref{eq:m24:dims} divides $|\Mtf|$, as the degree of an
irreducible representation must (Frobenius; standard, not in this book's library); a
symbolic check confirms it for all $26$.
\end{example}

\begin{pitfallbox}[What Burnside's check does and does not prove]
Passing \eqref{eq:m24:burnside} is strong evidence that a transcribed table is correct,
because a single wrong, duplicated or missing entry changes the sum by a large amount that
is unlikely to be compensated. It is \emph{not} a proof that the $26$ numbers are the
degrees of $\Mtf$: infinitely many other $26$-tuples of positive integers have the same sum
of squares. And the check is only as good as the value of $|G|$ it is compared with, which
must come from an independent source (here \eqref{eq:m24:order}, EOT eq.~(A.1)).
\end{pitfallbox}

\subsection*{The moonshine numbers, briefly}
EOT's coefficients $A_n$ of \cref{ch:ellgenus} are $45, 231, 770, 2277, 5796, 13915, 30843,
\dots$ \source{EOT eq.~(1.12), ll.~197--203}; the first five are single entries of
\eqref{eq:m24:dims}, and the next two decompose as $A_6=3520+10395$ and
$A_7=10395+5796+5544+5313+2024+1771$ \source{EOT eqs.~(1.14)--(1.15), ll.~228--239}. GHV
absorb a factor $2$ and write $90=45+\overline{45}$, $462=231+\overline{231}$,
$1440=770+\overline{770}$, $4554=2\cdot2277$, $11592=2\cdot5796$, $27830=2\cdot3520+2\cdot10395$,
$61686=2(1771+2024+5313+5796+5544+10395)$ \source{GHV eqs.~(3.3)--(3.4), ll.~555--575}, so
that each $2A_n$ is the dimension of a \emph{real} representation. \Cref{ch:moonshine} takes
the story from here; for now note only that the decompositions for $n\ge6$ are not unique
\source{EOT §1, ll.~237--239}, which is why a stronger test (the twining genera of GHV) was
needed.

%=====================================================================================
\section{\texorpdfstring{$\Mtf$}{M24} and the geometry of K3}\label{sec:m24:k3}

Let $S$ be a K3 surface and $G$ a finite group of automorphisms of $S$ that preserve the
holomorphic two-form (\emph{symplectic} automorphisms, \cref{ch:k3}). Write
$\Lambda=H^*(S,\Z)$ for the Mukai lattice of \cref{ch:mukai}, even, self-dual, of signature
$(4,20)$; let $\Lambda^G$ be the sublattice fixed by $G$ and $\Lambda_G$ its orthogonal
complement. EOT summarize Mukai's argument as follows \source{EOT App.~B, ll.~984--1001}.
\begin{enumerate}[leftmargin=*,itemsep=2pt]
\item $\Lambda_G$ lies in the primitive part of $H^{1,1}$ and is negative definite.
\item By Nikulin's embedding results, $\Lambda_G$ embeds in the Niemeier lattice
 $N=N_{24}$ of \cref{sec:m24:golay}, hence $G$ acts on $N$ and $G\subseteq\Mtf$.
\item $G$ fixes $H^0$, $H^4$, $H^{2,0}$, $H^{0,2}$ and the K\"ahler class, so
 $\rk\Lambda^G\ge5$ and $\rk\Lambda_G\le19$. Transported to the action on $24$ points, $G$
 has at least $5$ orbits.
\item Conversely, a subgroup of $\Mtf$ with at least $5$ orbits on the $24$ points acts on
 $H^{1,1}$ of some K3 surface, by the global Torelli theorem (\cref{ch:k3lattice}).
\end{enumerate}
This is Mukai's theorem\index{Mukai's theorem} (Tier~L: \source{EOT App.~B, ll.~955--958,
1054--1057} cite Mukai 1988 and Kondo 1998; the proofs are not in this book's library).
In particular $G$ is never all of $\Mtf$, which acts transitively (one orbit). EOT's example
is the Fermat quartic $X^4+Y^4+Z^4+W^4=0$ in $\mathbb{CP}^3$, whose symplectic symmetry group
$(\Z_4)^2\rtimes S_4$ has order $16\cdot24=384=2^7\cdot3$ and decomposes the $24$ points into
five orbits of lengths $1,1,2,4,16$ \source{EOT App.~B, ll.~1003--1006}; note $1+1+2+4+16=24$
and $384\mid|\Mtf|$.

\begin{physicsbox}[Geometric symmetry versus moonshine symmetry]
Mukai's theorem gives a \emph{different} subgroup of $\Mtf$ at each special point of the K3
moduli space, never the whole group. The moonshine observation of \cref{ch:moonshine}
concerns the elliptic genus, which is constant over the moduli space; EOT ask whether the
various geometric subgroups are ``enhanced to $\Mtf$ over the whole of moduli space'' at the
level of the elliptic genus \source{EOT §1, ll.~268--275}. In this book that enhancement is
a Tier~L statement of the moonshine literature where cited, and Tier~C where this
programme extends it.
\end{physicsbox}

%=====================================================================================
\section{How a wrong table was caught}\label{sec:m24:story}

The project's table of $\Mtf$ representation dimensions, \lean{M24RepDim}, was entered by
hand from EOT eq.~(A.3) and was ``verified'' for months, in the sense that the library
compiled and its one theorem about the table, a lookup of the second entry ($23$), passed.
The table was nevertheless wrong: it listed $10395$ and $483$ twice and omitted the second
$990$ and the third $1035$ \source{MathieuM24.lean, ll.~65--67; LL.md §S2.6, ll.~63--67}.
Such a table has the right length ($26$), the right first entries, and the right multiset
of \emph{distinct} values; a lookup theorem cannot see the difference. Burnside's identity
can. With the two wrong substitutions the sum of squares changes by
\begin{equation}\label{eq:m24:wrongsum}
 10395^2+483^2-990^2-1035^2=108\,056\,025+233\,289-980\,100-1\,071\,225=106\,237\,989,
\end{equation}
so the wrong table sums to $244\,823\,040+106\,237\,989=351\,061\,029\ne|\Mtf|$, and the
theorem \lean{M24RepDim\_sum\_sq} of \cref{sec:m24:lean} fails to compile against it. Once
the table was corrected on 2026-09-16, the theorem closed by unfolding the $26$-term sum
and normalizing numerals \source{MathieuM24.lean, ll.~39--45, 117--121}.

\begin{historybox}[The lesson, as the project recorded it]
``Guard every data table with a theorem that would break if the table were wrong.
Stream~1's \lean{M24RepDim} table had been `verified' for months and was wrong (two entries
duplicated, two missing), because its only theorem checked a single entry. Burnside's
$\sum\dim^2=|G|$ caught it in one line. For every hand-entered table, add the cheapest
global consistency theorem available (sum of squares, total count, known product,
symmetry).'' \source{LL.md §S2.6, ll.~63--67}
\end{historybox}

The general point is worth restating for a student of formal methods. A proof assistant
checks that a proof proves its statement; it does not check that the statement is the one
you meant, nor that the data in a definition were copied correctly. Global invariants that
the data must satisfy, sums, counts, factorizations, symmetries, are cheap to state, cheap to
prove by \lean{norm\_num} or \lean{decide}, and catch exactly the errors a lookup cannot.
Notice also the second consistency theorem in \lean{MathieuTower.lean}: the tower product
$24\cdot23\cdot22\cdot20160$ is compared not with a fresh literal but with the
\emph{already-certified} factorization \lean{M24\_order}, so that the two files cannot drift
apart silently \source{MathieuTower.lean, ll.~97--110}.

%=====================================================================================
\section{What the machine has checked}\label{sec:m24:lean}

\begin{sloppypar}
Five files touch $\Mtf$ or the Golay code. Two of them are worth reading:
\lean{MathieuM24.lean} and \lean{MathieuTower.lean}, both in the library
\lean{StringTheoryFormalization}. Three others, in the Mathlib-free libraries
\lean{Lean5Corpus} and \lean{StringTheoryFoundation}, are arithmetic shadows whose
docstrings promise much more than their statements deliver, and we describe them for that
reason. According to the
project's axiom audit \source{docs/VERIFIED\_FOUNDATION.md, ll.~7--16}, every theorem
below depends only on the axioms \lean{propext}, \lean{Classical.choice} and
\lean{Quot.sound}. \emph{No group, group action, code or representation is constructed in any of
these files}: they contain natural-number tables and equalities between numerals
\source{MathieuM24.lean, ll.~21--24, 110--114; MathieuTower.lean, ll.~43--46}.
\end{sloppypar}

\begin{leanbox}[The dimension table, the order and Burnside (\lean{MathieuM24.lean})]
\begin{leancode}
def M24RepDim : Fin 26 → ℕ
  | ⟨0, _⟩ => 1      | ⟨1, _⟩ => 23     | ⟨2, _⟩ => 45     | ⟨3, _⟩ => 45
  | ⟨4, _⟩ => 231    | ⟨5, _⟩ => 231    | ⟨6, _⟩ => 252    | ⟨7, _⟩ => 253
  | ⟨8, _⟩ => 483    | ⟨9, _⟩ => 770    | ⟨10, _⟩ => 770   | ⟨11, _⟩ => 990
  | ⟨12, _⟩ => 990   | ⟨13, _⟩ => 1035  | ⟨14, _⟩ => 1035  | ⟨15, _⟩ => 1035
  | ⟨16, _⟩ => 1265  | ⟨17, _⟩ => 1771  | ⟨18, _⟩ => 2024  | ⟨19, _⟩ => 2277
  | ⟨20, _⟩ => 3312  | ⟨21, _⟩ => 3520  | ⟨22, _⟩ => 5313  | ⟨23, _⟩ => 5796
  | ⟨24, _⟩ => 5544  | ⟨25, _⟩ => 10395
  | ⟨n + 26, h⟩ => absurd h (by omega)

theorem M24_first_coefficient :
    M24RepDim ⟨1, by norm_num⟩ = 23 := by rfl

theorem M24_order : (244823040 : ℕ) = 2^10 * 3^3 * 5 * 7 * 11 * 23 := by norm_num

theorem M24RepDim_sum_sq : (∑ i : Fin 26, M24RepDim i ^ 2) = 244823040 := by
  simp [Fin.sum_univ_succ, M24RepDim]
\end{leancode}
(The $26$ cases of \lean{M24RepDim} are written one per line in the source; they are
re-wrapped here.) \emph{Reading.} \lean{M24RepDim} is a function from the $26$-element type
\lean{Fin 26} to $\N$, a literal transcription of \eqref{eq:m24:dims} in EOT's order. The
final case is unreachable (\lean{n + 26 < 26} is refuted by \lean{omega}), so a missing entry
would be a compile error rather than a silent default \source{MathieuM24.lean, ll.~95--98}.
\lean{M24\_order} is the numeral identity \eqref{eq:m24:order}; \lean{M24RepDim\_sum\_sq} is
the computation of \cref{exa:m24:burnside}; \lean{M24\_first\_coefficient} is a single lookup
kept as a regression guard. \emph{Proof idea.} \lean{norm\_num} evaluates both sides of
\lean{M24\_order} to the same numeral. For the sum, \lean{Fin.sum\_univ\_succ} peels the
sum over \lean{Fin 26} into $26$ explicit terms, each \lean{M24RepDim} application reduces
to its literal, and \lean{simp}'s arithmetic normalizes the result. \emph{Scope.} Burnside's
theorem itself (\cref{prop:m24:burnside}) is not formalized; that $244\,823\,040$ is the
order of a group is not formalized; the identification of the table with $\Mtf$ is Tier~L.
\end{leanbox}

\begin{leanbox}[The stabilizer tower (\lean{MathieuTower.lean})]
\begin{leancode}
def orderM23 : ℕ := 10200960
def orderM22 : ℕ := 443520
def orderM21 : ℕ := 20160

theorem M24_stabilizer_index : 244823040 = 24 * orderM23 := by ...
theorem M23_stabilizer_index : orderM23 = 23 * orderM22 := by ...
theorem M22_stabilizer_index : orderM22 = 22 * orderM21 := by ...
theorem M21_order_factorization : orderM21 = 2 ^ 6 * 3 ^ 2 * 5 * 7 := by ...

theorem mathieu_tower_consistent :
    24 * (23 * (22 * orderM21)) = 244823040 := by
  rw [← M22_stabilizer_index, ← M23_stabilizer_index, ← M24_stabilizer_index]

theorem mathieu_tower_matches_certified_order :
    24 * (23 * (22 * orderM21)) = 2 ^ 10 * 3 ^ 3 * 5 * 7 * 11 * 23 := by
  rw [mathieu_tower_consistent, M24_order]
\end{leancode}
\emph{Reading.} The three orders \eqref{eq:m24:M23}--\eqref{eq:m24:M21} are hard-coded
constants, and the theorems are exactly the index relations \eqref{eq:m24:indices}, the
factorization \eqref{eq:m24:M21}, and their composition. \emph{Proof idea.} Each index
relation is \lean{unfold} followed by \lean{norm\_num}. The two composition theorems are
pure rewriting: \lean{mathieu\_tower\_consistent} substitutes the three index equations from
the inside out until both sides are the same numeral, and the last theorem chains it with
the imported \lean{M24\_order} \source{MathieuTower.lean, ll.~48--53}. \emph{Scope.} As the
docstring itself says, this ``is \emph{not} a formalization of the Mathieu group, its
Golay-code action, transitivity, or the orbit--stabilizer theorem itself''
\source{MathieuTower.lean, ll.~43--46}. The ATLAS orders are Tier~L and are not in this
book's library of sources.
\end{leanbox}

\begin{leanbox}[Golay-code numerals (\lean{Problem8\_GolayHolography.lean},
\lean{TensorNetworkBridge.lean})]
\begin{leancode}
def golay_length : Nat := 24
def golay_dimension : Nat := 12
def golay_distance : Nat := 8
def golay_codeword_count : Nat := 2 ^ 12
def error_correction_radius (d : Nat) : Nat := (d - 1) / 2

theorem golay_codeword_count_eval : golay_codeword_count = 4096 := by rfl
theorem golay_error_radius_equals_three :
    error_correction_radius golay_distance = 3 := by rfl
theorem golay_self_dual_dimension :
    golay_length - golay_dimension = golay_dimension := by rfl
theorem nonzero_codeword_weight_positive (w : Nat)
    (h : is_valid_nonzero_codeword_weight w) : w > 0 := by ...
\end{leancode}
and, in the second file,
\begin{leancode}
structure GolayCodeParameters where
  length : Nat := 24
  dimension : Nat := 12
  distance : Nat := 8

theorem golay_length_matches_k3_euler :
    (defaultGolay.length : Int) = eulerChar4D bettiK3 := by decide
theorem golay_code_rate_is_half :
    defaultGolay.dimension * 2 = defaultGolay.length := by rfl
\end{leancode}
\emph{Reading.} These files define the three parameters $[24,12,8]$ as natural numbers and
prove $2^{12}=4096$, $\lfloor(8-1)/2\rfloor=3$, $24-12=12$, and $12\cdot2=24$. The one
statement with any cross-file content is \lean{golay\_length\_matches\_k3\_euler}: the
literal $24$ equals the Euler characteristic computed from the K3 Betti numbers
\lean{bettiK3} $=(1,0,22,0,1)$ by \lean{eulerChar4D} in the project's \lean{Topology}
module, which \lean{euler\_char\_K3} evaluates to $24$. \lean{nonzero\_codeword\_weight\_positive}
says that a natural number $w\ge8$ is $>0$. \emph{Scope.} No code, no codeword, no
Hamming weight and no automorphism appears in either file; the predicate
\lean{is\_valid\_nonzero\_codeword\_weight w} is literally \lean{w ≥ 8}. The docstring of
\lean{Problem8} asserts that these theorems ``prove that any bulk local operator on
$K3\times T^2$ can be perfectly reconstructed'' after erasure of three boundary regions; they
prove nothing of the sort, and the holographic reading is Tier~C at best. The second file
carries a scope note added after review saying the same of its tensor-network narrative
\source{TensorNetworkBridge.lean, ll.~39--46}.
\end{leanbox}

\begin{leanbox}[Numerology about $|\Mtf|$ (\lean{Problem2\_MathieuFrobeniusRigidity.lean})]
\begin{leancode}
def bps_conductor : Nat := 27720
def order_M24 : Nat := 244823040

theorem bps_conductor_prime_factorization :
    (2 ^ 3) * (3 ^ 2) * 5 * 7 * 11 = bps_conductor := by rfl
theorem m24_conductor_quotient :
    order_M24 / bps_conductor = 8832 ∧ order_M24 % bps_conductor = 0 := by decide
theorem hasse_trace_integer_bound (a : Int) (h : a * a ≤ 52) :
    a * a ≤ 52 := h
\end{leancode}
\emph{Reading.} $27720=2^3\cdot3^2\cdot5\cdot7\cdot11$ and $244\,823\,040=27720\cdot8832$
(with $8832=2^7\cdot3\cdot23$; symbolic check). Both are true numeral facts. The number
$27720$ is this programme's ``BPS character lock'' ($=462\cdot60=4\cdot90\cdot77$), a
Tier~C numerological proposal with no source in this book's library; the divisibility by
$|\Mtf|$ is arithmetic and is also re-proved in
\lean{DualScaleValidation.UseCase2.m24\_order\_divisible\_by\_bps\_lock}. The theorem
\lean{hasse\_trace\_integer\_bound} deserves a moment's attention: its conclusion is its
hypothesis, and the proof is the hypothesis itself. It proves nothing, and the ``Hasse--Weil
bound'' of the docstring is not formalized. A reader auditing a library should look for
exactly this pattern.
\end{leanbox}

\begin{sloppypar}
Finally, the dimension table is a \emph{hub}. Four theorems of the moonshine module
\lean{DualScaleStream2.Moonshine} (\lean{first\_five\_are\_irreps},
\lean{A6\_decomposition}, \lean{A7\_decomposition}, \lean{A6\_not\_irrep}) import
\lean{M24RepDim} by name and verify EOT's decompositions of $A_1,\dots,A_7$ against it
\source{MathieuM24.lean, ll.~47--52}. They belong to \cref{ch:moonshine}. Here we only
note that they inherit the correctness of the table: had \lean{M24RepDim} still contained
a second $483$ in place of the second $990$, \lean{first\_five\_are\_irreps} would still
have compiled (its five witnesses are at indices $2,4,9,19,23$), which is another
illustration that lookups do not guard tables.
\end{sloppypar}

\begin{center}\small
\begin{tabular}{@{}>{\raggedright\arraybackslash}p{0.43\textwidth}c
  >{\raggedright\arraybackslash}p{0.45\textwidth}@{}}
\toprule
Claim & Tier & Evidence\\
\midrule
$244\,823\,040=2^{10}\cdot3^3\cdot5\cdot7\cdot11\cdot23$ & \TA & \lean{M24\_order}\\[2pt]
The $26$ entries of \eqref{eq:m24:dims} satisfy $\sum\dim^2=244\,823\,040$; entry $1$ is
 $23$ & \TA & \lean{M24RepDim\_sum\_sq}, \lean{M24\_first\_coefficient}\\[2pt]
Index relations \eqref{eq:m24:indices} for the ATLAS numerals; $20160=2^6\cdot3^2\cdot5\cdot7$;
 $24\cdot23\cdot22\cdot20160=|\Mtf|$ and equals the certified factorization & \TA &
 \lean{M24\_stabilizer\_index}, \lean{M23\_stabilizer\_index}, \lean{M22\_stabilizer\_index},
 \lean{M21\_order\_factorization}, \lean{mathieu\_tower\_consistent},
 \lean{mathieu\_tower\_matches\_certified\_order}\\[2pt]
$2^{12}=4096$; $\lfloor7/2\rfloor=3$; $24-12=12$; $12\cdot2=24$; $24=\chi(K3)$ from Betti
 numbers & \TA & \lean{golay\_codeword\_count\_eval}, \lean{golay\_error\_radius\_equals\_three},
 \lean{golay\_self\_dual\_dimension}, \lean{golay\_code\_rate\_is\_half},
 \lean{golay\_length\_matches\_k3\_euler}\\[2pt]
$27720=2^3\cdot3^2\cdot5\cdot7\cdot11$ divides $244\,823\,040$ with quotient $8832$ & \TA &
 \lean{bps\_conductor\_prime\_factorization}, \lean{m24\_conductor\_quotient}\\[2pt]
Lattice-from-code dictionary \eqref{eq:m24:pairing}--\eqref{eq:m24:roots}; $|PSL(3,4)|$;
 Burnside's identity; $5$-transitivity count \eqref{eq:m24:fivetrans} & --- & derived on the
 page\\[2pt]
Weight distribution \eqref{eq:m24:weights}; $[8,4,4]$ code has $14$ words of weight $4$;
 all degrees divide $|\Mtf|$ & --- & symbolic checks (scripts)\\[2pt]
$244\,823\,040$ is the order of $\Mtf$; $26$ classes and irreps; the dimensions
 \eqref{eq:m24:dims} and the conjugate pairs; uniqueness of $\mathcal G_{24}$;
 $\Mtf=\operatorname{Aut}(\mathcal G_{24})$ & \TL & EOT App.~A, ll.~331--358; App.~B,
 ll.~959--979\\[2pt]
The $26$ class names; power-map and $M_{23}$ facts & \TL & GHV ll.~694--697, 808--813,
 1916--1930\\[2pt]
$5$-transitivity; ATLAS orders of $M_{23},M_{22},M_{21}$; $M_{21}\cong PSL(3,4)$; Steiner
 system; degrees divide the order & \TL & standard; not in this book's library\\[2pt]
Mukai's theorem; Fermat quartic with $384$ symmetries and orbits $1,1,2,4,16$ & \TL &
 EOT App.~B, ll.~984--1006\\[2pt]
Golay code as holographic error-correcting code on $\KT$; $27720$ as a ``BPS lock'' &
 \TC & docstrings of \lean{Problem8}, \lean{Problem2}; not supported by any theorem\\
\bottomrule
\end{tabular}
\end{center}

%=====================================================================================
\begin{summarybox}
\begin{itemize}[leftmargin=*,itemsep=2pt]
\item $\Mtf$ is the automorphism group of the extended binary Golay code $\mathcal G_{24}$,
 the unique doubly even self-dual $[24,12,8]$ code; equivalently, the permutation
 symmetries of the Niemeier lattice $A_1^{24}\subset N\subset(A_1^*)^{24}$. The dictionary
 code $\leftrightarrow$ lattice is: doubly even $\leftrightarrow$ even, self-dual
 $\leftrightarrow$ unimodular, no weight-$4$ words $\leftrightarrow$ no extra roots. With
 $n=8$ the same recipe gives $E_8$ with $16+14\cdot16=240$ roots.
\item $|\Mtf|=2^{10}\cdot3^3\cdot5\cdot7\cdot11\cdot23=244\,823\,040$. The action on $24$
 points is $5$-transitive, and fixing points gives $M_{23}$ ($10\,200\,960$), $M_{22}$
 ($443\,520$), $M_{21}\cong PSL(3,4)$ ($20\,160$), with indices $24,23,22$; five points have
 a stabilizer of order $48$.
\item There are $26$ conjugacy classes and $26$ irreducible representations, of dimensions
 \eqref{eq:m24:dims}, five complex-conjugate pairs among them. Burnside's
 $\sum\dim^2=|G|$ holds, and it caught a wrong table in this project.
\item Mukai: finite symplectic automorphism groups of K3 surfaces are the subgroups of
 $\Mtf$ with at least five orbits on $24$ points; never $\Mtf$ itself.
\item Lean has checked the arithmetic: the order factorization, the Burnside sum for the
 transcribed table, the tower indices, and the Golay numerals. No group, code or
 representation is constructed; the identification with $\Mtf$ is Tier~L.
\end{itemize}
\end{summarybox}

%=====================================================================================
\section*{Exercises}
\addcontentsline{toc}{section}{Exercises}

\begin{exercise}[$\star$ Lagrange at work]\label{ex:m24:lagrange}
Using only \eqref{eq:m24:order}, decide which of the following can be the order of an
element of $\Mtf$: $9$, $13$, $16$, $22$, $25$, $33$, $46$, $69$. Then compare with the list
of element orders read off \eqref{eq:m24:classes}: which orders are permitted by Lagrange
but do not occur? (This shows that Lagrange gives necessary conditions only.)
\end{exercise}

\begin{exercise}[$\star$ The Hamming code by hand]\label{ex:m24:hamming}
List all $16$ codewords of the code $C_8$ of \cref{exa:m24:hamming} and verify the weight
distribution $1,14,1$. Check that the complement $c\mapsto c+(1,\dots,1)$ maps weight-$4$
words to weight-$4$ words, and use this to explain why the distribution is symmetric.
\end{exercise}

\begin{exercise}[$\star\star$ Construction A in general]\label{ex:m24:constrA}
Let $C\subseteq\mathbb F_2^n$ be any binary code and $N_C$ the lattice \eqref{eq:m24:sandwich}.
(a) Show that $\operatorname{disc}N_C=2^{\,n-2\dim C}$. (b) Show that $N_C$ is integral
iff $C\subseteq C^\perp$, and even iff $C$ is doubly even. (c) Show that the minimum norm
of a nonzero vector of $N_C$ is $\min(2,\,d/2)$ where $d$ is the minimum distance, and
conclude that for $\mathcal G_{24}$ the minimum norm is $2$ but no root comes from a
nonzero codeword.
\end{exercise}

\begin{exercise}[$\star\star$ Burnside for a small group]\label{ex:m24:S4}
The symmetric group $S_4$ has five conjugacy classes (cycle types $1^4$, $2\,1^2$, $2^2$,
$3\,1$, $4$) and irreducible representations of dimensions $1,1,2,3,3$. Verify
\eqref{eq:m24:burnside}. Then show that no list of five positive integers other than a
permutation of $1,1,2,3,3$ has sum of squares $24$ \emph{and} contains $1$ twice (the
trivial and sign representations); this is a toy version of the ``strong evidence, not
proof'' remark of \cref{sec:m24:chars}.
\end{exercise}

\begin{exercise}[$\star\star$ Five points]\label{ex:m24:fivepoints}
(a) Verify \eqref{eq:m24:fivetrans}. (b) A group acting $k$-transitively on $n$ points is
\emph{sharply} $k$-transitive if the pointwise stabilizer of $k$ points is trivial. Show that
$\Mtf$ is not sharply $5$-transitive but that, if it were sharply $t$-transitive for some
$t$, its order would be $24\cdot23\cdots(25-t)$; check that no $t$ works. (c) $M_{12}$, of
order $95\,040$ (standard value; not in this book's library), is sharply $5$-transitive on
$12$ points: verify $12\cdot11\cdot10\cdot9\cdot8=95\,040$.
\end{exercise}

\begin{exercise}[$\star\star$ Orbits of the Fermat quartic group]\label{ex:m24:fermat}
The group $(\Z_4)^2\rtimes S_4$ of \cref{sec:m24:k3} has order $384$ and orbits of lengths
$1,1,2,4,16$ on the $24$ points. For each orbit, use orbit--stabilizer to compute the order
of the stabilizer of a point in it, and check that all five are divisors of $384$. Which
orbit has the smallest stabilizer, and what does Mukai's ``at least five orbits'' condition
say about $\rk\Lambda^G$ here?
\end{exercise}

\begin{exercise}[$\star\star$ Lean: guard the tower table]\label{ex:m24:leantower}
In the style of \lean{M21\_order\_factorization}, state and prove in Lean the factorizations
$\lean{orderM23} = 2^7\cdot3^2\cdot5\cdot7\cdot11\cdot23$ and
$\lean{orderM22}=2^7\cdot3^2\cdot5\cdot7\cdot11$ from \eqref{eq:m24:M23}--\eqref{eq:m24:M22}
(\lean{unfold} then \lean{norm\_num}). Explain why these two theorems would have caught a
mistyped \lean{orderM22} that the three index theorems alone might not (hint: the index
theorems constrain the three constants only up to a common factor if all three are
mistyped consistently).
\end{exercise}

\begin{exercise}[$\star\star\star$ Lean: degrees divide the order]\label{ex:m24:leandiv}
State the theorem \lean{∀ i : Fin 26, 244823040 \% M24RepDim i = 0} and prove it by
\lean{decide} or by \lean{fin\_cases i <;> norm\_num [M24RepDim]}. Then explain in words why
this theorem, like \lean{M24RepDim\_sum\_sq}, is a statement about $26$ numerals and not
about representations of $\Mtf$, and why it is nevertheless a useful guard (which
transcription errors would it catch that the Burnside sum might not, and vice versa?).
\end{exercise}

\begin{exercise}[$\star\star\star$ The wrong table, quantitatively]\label{ex:m24:wrongtable}
(a) Verify \eqref{eq:m24:wrongsum}. (b) Suppose a single entry $d$ of \eqref{eq:m24:dims}
is replaced by another entry $d'$ of the same table (a duplication error). Show that the
Burnside sum is unchanged iff $d'=d$, so \emph{every} single duplication is caught. (c) Find
the smallest change to the table that Burnside would \emph{not} catch: two entries
$d_1,d_2$ replaced by $d_1',d_2'$ with $d_1^2+d_2^2=d_1'^2+d_2'^2$, all four being
positive integers dividing $|\Mtf|$ (a script is acceptable; you may find none, in which
case say what you searched).
\end{exercise}

%=====================================================================================
\section*{Further reading}
\addcontentsline{toc}{section}{Further reading}

\begin{itemize}[leftmargin=*,itemsep=2pt]
\item \source{EOT App.~A, ll.~329--358} is the source of the order, the count of classes and
 representations, the dimension list (A.3) and the conjugate pairs (A.4); the full character
 table follows at ll.~360--925 but is garbled by text extraction and was not used in this
 chapter. \source{EOT App.~B, ll.~955--1007} contains the Golay-code definition of $\Mtf$
 via $A_1^{24}$ and the summary of Mukai's theorem, with references to Mukai 1988 and
 Kondo 1998 at ll.~1054--1057.
\item \source{GHV §3, ll.~500--575} restates the decompositions with the factor $2$
 absorbed; \source{GHV Table~1, ll.~1916--1930} gives the class names and the power maps
 used in \cref{sec:m24:chars}; \source{GHV §3.1, ll.~808--813} identifies which classes meet
 $M_{23}$.
\item Lean sources: \lean{MathieuM24.lean} (directory \lean{StringDynamics/}) and
 \lean{MathieuTower.lean} (directory \lean{UseCases/}), both in the library
 \lean{StringTheoryFormalization}. The scope discussion at ll.~8--57 of the first file is a
 model of how a docstring should say what is and is not proved. \source{LL.md §S2.6,
 ll.~63--67} records the lesson of \cref{sec:m24:story}.
\item The classification of finite simple groups, the ATLAS of finite groups, the Steiner
 system $S(5,8,24)$ and the uniqueness of the Golay code are standard material of finite
 group theory and coding theory; none of it is in this book's library, and a student who
 wants the proofs should consult a group-theory text.
\end{itemize}
